\documentclass[11pt, a4paper, oneside]{article}
\usepackage[utf8]{inputenc}
\usepackage[T1]{fontenc}
\usepackage{geometry}
\usepackage{amsmath, amssymb}
\usepackage{fancyhdr}
\usepackage{hyperref}
\usepackage{booktabs}
\usepackage{setspace}

\geometry{margin=2cm}
\pagestyle{fancy}
\fancyhf{}
\rhead{Exposé}
\cfoot{\thepage}
\setlength{\headheight}{15pt}
\hypersetup{colorlinks=true, linkcolor=blue, urlcolor=cyan}
\setlength{\parindent}{0pt}
\setlength{\parskip}{0.5ex}
\setlength{\intextsep}{0.5em}

\begin{document}
\pagestyle{empty}

\vspace*{-2cm}
\begin{center}
{\huge\textbf{Exposé}}\\[2mm]
{Masterarbeit / Bachelorarbeit}\\\\[4mm]
\textbf{\Large Title Here}

\vspace{4mm}
\begin{tabular}{@{}p{0.55\textwidth}p{0.35\textwidth}@{}}
\textbf{Student:} Name & \textbf{Supervisor:} \\
University & Prof.\ Name \\
& Department, University \\
& \\
\today & \\
\end{tabular}
\end{center}

\newpage

\section{Forschungsfrage und Motivation}

\subsection{Hintergrund und Problemstellung}
Replace with your problem description. This section should establish:
\begin{itemize}
    \item The broader context (what field/domain?)
    \item The specific problem (what goes wrong today?)
    \item Why it matters (consequences of not solving it?)
\end{itemize}

\subsection{Forschungsfrage}
\begin{quote}
\textit{Your research question here.}
\end{quote}

\subsection{Ziele}
\begin{itemize}
    \item \textbf{Primär:} Main quantitative goal.
    \item \textbf{Sekundär:} Model comparison or methodological goal.
    \item \textbf{Tertiär:} Analysis or diagnostic goal.
\end{itemize}

\section{Stand der Forschung}

\subsection{Subsection 1}
Review of relevant work in this area.

\subsection{Subsection 2}
Gap analysis -- what existing approaches miss.

\subsection{Subsection 3}
How your approach differs and why that matters.

\section{Methodik}

\subsection{Workload/Experimental Model}
Describe your data, simulation, or experimental setup.

\subsection{ML-Modelle (or Methods)}
List your methods/models with key hyperparameters.

\subsection{Simulationsaufbau und Bewertung}
Describe the comparison setup and evaluation metrics.

\section{Erwartete Ergebnisse, Abgrenzung und Planung}

\subsection{Erwartete Ergebnisse}
List your expected findings with quantitative targets.

\subsection{Abgrenzung}
What is excluded from the scope.

\subsection{Arbeitsplanung}
\begin{table}[h!]
\centering\small
\begin{tabular}{@{}llc@{}}
\toprule
\textbf{Phase} & \textbf{Beschreibung} & \textbf{Dauer} \\
\midrule
1. Literaturstudium & Overview & X Wchn. \\
2. Methodik & Implementation & X Wchn. \\
3. Modellentwicklung & Training and validation & X Wchn. \\
4. Auswertung & Analysis & X Wchn. \\
5. Schriftfassung & Writing & X Wchn. \\
6. Abschluss & Review and defense & X Wchn. \\
\midrule
& \textbf{Gesamtdauer} & \textbf{ca.\ XX Wchn.} \\
\bottomrule
\end{tabular}
\end{table}

\subsection{Risikoabschätzung}
\begin{itemize}
    \item \textbf{Gering:} Low-risk items.
    \item \textbf{Mittel:} Medium-risk items + mitigation.
\end{itemize}

\newpage
\bibliographystyle{ieeetr}
\begin{thebibliography}{9}

\bibitem{placeholder}
Author, A. (Year).
\textit{Title}.
Journal/Conference.

\end{thebibliography}

\end{document}
